% ICME 2027 skeleton -- The Price of a Model-Free 3DGS Bitstream
% Build: pdflatex main && bibtex main && pdflatex main && pdflatex main
\documentclass[conference]{IEEEtran}
\usepackage{graphicx}
\usepackage{booktabs}
\usepackage{amsmath}
\usepackage{amssymb}
\usepackage{url}
\usepackage{xcolor}
\usepackage[pagebackref,breaklinks,colorlinks,allcolors=blue]{hyperref}

\newcommand{\dcca}{DCCA-GS}
\newcommand{\todo}[1]{\textcolor{red}{[TODO: #1]}}

\begin{document}

\title{The Price of a Model-Free 3DGS Bitstream:\\
A Scene-Dependent Measurement and Decision Rule}

\author{\IEEEauthorblockN{Anonymous Authors}
\IEEEauthorblockA{Anonymous Institution\\
Email: anonymous@example.com}}

\maketitle

\begin{abstract}
Anchor-based 3D Gaussian Splatting (3DGS) codecs compress a scene to a few megabytes, but their entropy model and decoder parameters must travel with the bitstream. A recent line removes this dependency: quantization steps and probability conditioning are recomputed from decoded bytes, and only kilobyte-scale static tables are carried inline. The benefit is deployment independence; the cost is unknown. We measure this cost on seven public scenes under a single protocol. At matched quality, the model-free stream is 0.2--0.3\,dB behind a learned single-rate codec on garden and train, but 1.3\,dB behind on truck, while its full-precision file is 42--48\% larger. We further profile decoding: ten semantic chunks take 14.3\,s on a server CPU, while steady-state rendering takes 6--7\,s per tier. These numbers yield a simple decision rule: a model-free progressive bitstream pays off when model versioning or bandwidth uncertainty dominates, and the scene is in the codec's comparable regime; otherwise a learned single-rate codec is smaller and better. We release the protocol, per-scene numbers, and the decision table.
\end{abstract}

\begin{IEEEkeywords}
3D Gaussian Splatting, compression, progressive bitstream, model-free decoding, measurement
\end{IEEEkeywords}

\section{Introduction}
\label{sec:intro}
Anchor-based 3DGS codecs~\cite{lu2024,chen2024,chen2025} compress a scene to a few megabytes, but the entropy model and decoder parameters are part of the bitstream. Recent progressive codecs~\cite{chen2026pcgs,tang2026progs,revilla2026scar,disario2025gode} improve rate--quality, yet all carry learned side information ahead of content. \dcca{} removes this dependency: quantization steps and probability conditioning are recomputed from decoded bytes, and the only priors are kilobyte-scale static tables. The benefit is deployment independence; the cost is a 42--48\% larger full-precision file.

This paper asks a narrower question than ``which codec is best'': \emph{when is the model-free premium worth paying?} We measure the premium on seven public scenes, decompose it into entropy-model loss and framing, profile decoding, and derive a decision rule.

\noindent\textbf{Contributions.}
\begin{itemize}
  \item We decompose the model-free premium into three parts: the learned-context entropy gain, the static-table/framing cost, and the header itself. The dominant term is the lost entropy model, not the header.
  \item We report a scene-dependent base comparison on seven public scenes: garden and train are within 0.3\,dB, while truck is 1.3\,dB behind at a similar size. A single average hides this reversal.
  \item We profile decoding and derive a decision rule for when a model-free progressive bitstream is preferable to a learned single-rate file.
\end{itemize}

\section{Related Work}
\label{sec:related}
\paragraph{Single-rate anchor compression}
Scaffold-GS~\cite{lu2024} introduces anchors; HAC/HAC++~\cite{chen2024,chen2025} add hash-grid context and adaptive quantization; ContextGS and CAT-3DGS improve context modeling. Their entropy model and decoder are part of the stream.

\paragraph{Progressive 3DGS}
PCGS~\cite{chen2026pcgs} uses progressive masking and quantization; ProGS~\cite{tang2026progs} uses parent-closed octree prefixes; SCAR-GS~\cite{revilla2026scar} uses RVQ and autoregressive context; GoDe~\cite{disario2025gode} uses gradient-based layer assignment and Zstd dictionaries. All carry learned side information; their headers are part of the rate.

\paragraph{Model-free and 2D precedents}
OIC-GS~\cite{oicgs2026} establishes zero-side-information context for 2D panoramas, but its header still carries MLP weights (32\% at the low-rate end). \dcca{} removes learned components from a 3D anchor stream. We use ``model-free'' rather than ``zero side information'' because static tables are still side information in the information-theoretic sense.

\section{Measurement Protocol}
\label{sec:protocol}
\paragraph{Scenes}
Seven public scenes: garden, bicycle, stump (Mip-NeRF360 outdoor), bonsai, room (indoor), train, truck (Tanks and Temples). All at the official factor/width, 30k iterations, seed 42. The DCCA evaluation factor equals the training factor; this avoids the resolution mismatch that inflated earlier gaps.

\paragraph{Comparisons}
We compare (i) at matched quality, (ii) at matched size by interpolation, and (iii) at matched decode budget. We report PSNR, SSIM, LPIPS, MB, encode/decode time, and TTFF. We do not report a single BD-rate when the PSNR intervals do not overlap.

\paragraph{Statistics}
Three scenes (garden, train, truck) use two seeds; the remaining four use one seed. We report per-scene values and ranges, not only the mean.

\section{Results}
\label{sec:results}

\subsection{The Base Is Scene-Dependent}
Table~\ref{tab:main} reports the seven-scene comparison. Garden and train are within 0.3\,dB; truck is 1.3\,dB behind at a similar size. The average would hide the reversal, so we report per scene.

\begin{table}[t]
\centering
\caption{Seven-scene comparison. DCCA uses $\lambda=5\times10^{-4}$ for the four new scenes; HAC++ uses its published low-rate point. Sizes in MB. \todo{fill after runs}}
\label{tab:main}
\begin{tabular}{lcccc}
\toprule
Scene & DCCA PSNR & DCCA MB & HAC++ PSNR & HAC++ MB \\
\midrule
garden & 26.82 & 9.46 & 27.01 & 8.99 \\
train  & \todo{} & \todo{} & 22.12 & 3.14 \\
truck  & \todo{} & \todo{} & 25.77 & 3.76 \\
bicycle& \todo{} & \todo{} & 25.12 & 13.89 \\
stump  & \todo{} & \todo{} & 26.53 & 10.71 \\
bonsai & \todo{} & \todo{} & 32.48 & 5.29 \\
room   & \todo{} & \todo{} & 31.46 & 3.51 \\
\bottomrule
\end{tabular}
\end{table}

\subsection{Where the Premium Goes}
Fig.~\ref{fig:premium} decomposes the full-precision file into symbol entropy, layer framing, geometry, and fixed payload. The learned-context gain of HAC++ is 33--40\%; the model-free file pays most of the 42--48\% premium for the missing entropy model. Header removal alone saves less than 3\,MB (PCGS 1.27--2.60\,MB; ProGS B\_ref 0.382\,MB; OIC-GS 32\% at low rate).

\begin{figure}[t]
\centering
\fbox{\parbox{0.9\linewidth}{\centering TODO: stacked bar of premium decomposition}}
\caption{Where the model-free premium goes.}
\label{fig:premium}
\end{figure}

\subsection{Decoding}
Ten semantic chunks take 14.331\,s in total; per-chunk decode is 0.335--2.224\,s. Steady-state rendering is 6.2--6.8\,s per tier; the first P3 render (152.5\,s) includes warm-up and is excluded. Decoding, not rendering, is the main steady-state cost.

\begin{figure}[t]
\centering
\fbox{\parbox{0.9\linewidth}{\centering TODO: per-chunk decode time + render line}}
\caption{Decode and render time.}
\label{fig:decode}
\end{figure}

\subsection{When It Pays}
Fig.~\ref{fig:decision} summarizes the decision rule. A model-free progressive stream is preferable when (i) model versions cannot be shipped or updated, (ii) bandwidth is unknown or heterogeneous so one file must serve multiple rates, (iii) the scene is in the comparable regime (garden/train-like), and (iv) the client can tolerate the decode cost or only needs a coarse tier. Otherwise a learned single-rate codec is smaller and better.

\begin{figure}[t]
\centering
\fbox{\parbox{0.9\linewidth}{\centering TODO: decision diagram (model version / bandwidth / scene / latency)}}
\caption{Decision rule.}
\label{fig:decision}
\end{figure}

\section{Limitations}
\label{sec:limitations}
Seven scenes are still a small sample; the truck reversal is measured but not explained; decoding is profiled on one server, not on an edge device; and the internal 1-78 UAV scene is not public. We do not claim SOTA compression, and we use ``model-free'' rather than ``zero side information''.

\section{Conclusion}
\label{sec:conclusion}
A model-free 3DGS bitstream is a deployment property, not a compression-rate property. Its cost can be measured (42--48\%), its benefit is conditional, and the decision rule above tells practitioners when to pay for it.

\bibliographystyle{IEEEtran}
\bibliography{references}

\end{document}
